\documentclass[12pt]{article}
\usepackage{amsmath,amssymb,amsfonts,mathrsfs}
\usepackage[margin=1in]{geometry}

\begin{document}

\begin{center}
{\Large\bfseries Chapter 10, Problem 21}
\end{center}

\noindent\textbf{Problem.} Consider all possible representations of the cubic group, $\mathscr{O}$, formed by taking the direct product of the various irreducible representations of $\mathscr{O}$. Find out how many times the irreducible representations of $\mathscr{O}$ occur in these product representations.

\bigskip
\noindent\textbf{Solution.}

The character table of $\mathscr{O}$ (from Eq.~(10.91)) is:

\[
\begin{array}{c|ccccc}
\mathscr{O} & \mathcal{C}_1 & \mathcal{C}_2 & \mathcal{C}_3 & \mathcal{C}_4 & \mathcal{C}_5 \\
& 1 & 6 & 3 & 8 & 6 \\
\hline
\Gamma_1 & 1 & 1 & 1 & 1 & 1 \\
\Gamma_2 & 1 & -1 & 1 & 1 & -1 \\
\Gamma_3 & 2 & 0 & 2 & -1 & 0 \\
\Gamma_4 & 3 & 1 & -1 & 0 & -1 \\
\Gamma_5 & 3 & -1 & -1 & 0 & 1
\end{array}
\]

The direct product $\Gamma_i \otimes \Gamma_j$ has character $\chi^{(i)}(g)\chi^{(j)}(g)$. We decompose each using $a_\nu = \frac{1}{24}\sum_k g_k \chi^{(\nu)*}_k \chi^{(i)}_k \chi^{(j)}_k$.

Since $\Gamma_1$ is the identity representation, $\Gamma_1 \otimes \Gamma_i = \Gamma_i$ for all $i$. Also $\Gamma_i \otimes \Gamma_j = \Gamma_j \otimes \Gamma_i$.

\medskip
\noindent$\boldsymbol{\Gamma_2 \otimes \Gamma_2}$: $\chi = (1, 1, 1, 1, 1) = \Gamma_1$.

\noindent$\boldsymbol{\Gamma_2 \otimes \Gamma_3}$: $\chi = (2, 0, 2, -1, 0) = \Gamma_3$.

\noindent$\boldsymbol{\Gamma_2 \otimes \Gamma_4}$: $\chi = (3, -1, -1, 0, 1) = \Gamma_5$.

\noindent$\boldsymbol{\Gamma_2 \otimes \Gamma_5}$: $\chi = (3, 1, -1, 0, -1) = \Gamma_4$.

\medskip
\noindent$\boldsymbol{\Gamma_3 \otimes \Gamma_3}$: $\chi = (4, 0, 4, 1, 0)$.

$a_1 = \frac{1}{24}[4 + 0 + 12 + 8 + 0] = 1$, $a_2 = \frac{1}{24}[4 + 0 + 12 + 8 + 0] = 1$, $a_3 = \frac{1}{24}[8 + 0 + 24 - 8 + 0] = 1$, $a_4 = \frac{1}{24}[12 + 0 - 12 + 0 + 0] = 0$, $a_5 = \frac{1}{24}[12 + 0 - 12 + 0 + 0] = 0$.

$\Gamma_3 \otimes \Gamma_3 = \Gamma_1 + \Gamma_2 + \Gamma_3$.

\medskip
\noindent$\boldsymbol{\Gamma_3 \otimes \Gamma_4}$: $\chi = (6, 0, -2, 0, 0)$.

$a_1 = \frac{1}{24}[6 + 0 - 6 + 0 + 0] = 0$, $a_2 = \frac{1}{24}[6 + 0 - 6 + 0 + 0] = 0$, $a_3 = \frac{1}{24}[12 + 0 - 12 + 0 + 0] = 0$, $a_4 = \frac{1}{24}[18 + 0 + 6 + 0 + 0] = 1$, $a_5 = \frac{1}{24}[18 + 0 + 6 + 0 + 0] = 1$.

$\Gamma_3 \otimes \Gamma_4 = \Gamma_4 + \Gamma_5$.

\medskip
\noindent$\boldsymbol{\Gamma_3 \otimes \Gamma_5}$: $\chi = (6, 0, -2, 0, 0)$.

Same characters as $\Gamma_3 \otimes \Gamma_4$! So $\Gamma_3 \otimes \Gamma_5 = \Gamma_4 + \Gamma_5$.

\medskip
\noindent$\boldsymbol{\Gamma_4 \otimes \Gamma_4}$: $\chi = (9, 1, 1, 0, 1)$.

$a_1 = \frac{1}{24}[9 + 6 + 3 + 0 + 6] = 1$, $a_2 = \frac{1}{24}[9 - 6 + 3 + 0 - 6] = 0$, $a_3 = \frac{1}{24}[18 + 0 + 6 + 0 + 0] = 1$, $a_4 = \frac{1}{24}[27 + 6 - 3 + 0 - 6] = 1$, $a_5 = \frac{1}{24}[27 - 6 - 3 + 0 + 6] = 1$.

$\Gamma_4 \otimes \Gamma_4 = \Gamma_1 + \Gamma_3 + \Gamma_4 + \Gamma_5$.

\medskip
\noindent$\boldsymbol{\Gamma_4 \otimes \Gamma_5}$: $\chi = (9, -1, 1, 0, -1)$.

$a_1 = \frac{1}{24}[9 - 6 + 3 + 0 - 6] = 0$, $a_2 = \frac{1}{24}[9 + 6 + 3 + 0 + 6] = 1$, $a_3 = \frac{1}{24}[18 + 0 + 6 + 0 + 0] = 1$, $a_4 = \frac{1}{24}[27 - 6 - 3 + 0 + 6] = 1$, $a_5 = \frac{1}{24}[27 + 6 - 3 + 0 - 6] = 1$.

$\Gamma_4 \otimes \Gamma_5 = \Gamma_2 + \Gamma_3 + \Gamma_4 + \Gamma_5$.

\medskip
\noindent$\boldsymbol{\Gamma_5 \otimes \Gamma_5}$: $\chi = (9, 1, 1, 0, 1)$.

Same as $\Gamma_4 \otimes \Gamma_4$: $\Gamma_5 \otimes \Gamma_5 = \Gamma_1 + \Gamma_3 + \Gamma_4 + \Gamma_5$.

\medskip
\noindent\textbf{Summary table of direct products:}

\[
\begin{array}{c|ccccc}
\otimes & \Gamma_1 & \Gamma_2 & \Gamma_3 & \Gamma_4 & \Gamma_5 \\
\hline
\Gamma_1 & \Gamma_1 & \Gamma_2 & \Gamma_3 & \Gamma_4 & \Gamma_5 \\
\Gamma_2 & & \Gamma_1 & \Gamma_3 & \Gamma_5 & \Gamma_4 \\
\Gamma_3 & & & \Gamma_1{+}\Gamma_2{+}\Gamma_3 & \Gamma_4{+}\Gamma_5 & \Gamma_4{+}\Gamma_5 \\
\Gamma_4 & & & & \Gamma_1{+}\Gamma_3{+}\Gamma_4{+}\Gamma_5 & \Gamma_2{+}\Gamma_3{+}\Gamma_4{+}\Gamma_5 \\
\Gamma_5 & & & & & \Gamma_1{+}\Gamma_3{+}\Gamma_4{+}\Gamma_5
\end{array}
\]

\end{document}
